% ECONOMICS_PROBLEM: {"id": "C31-282", "title": "Information and statistical power under interference", "jel_code": "C31", "source_id": "OP-0493"}
% FIRST_POSTED: 2026-10-09
% ATTEMPT_STATUS: unresolved
% ENTRY_KIND: research_attempt
\documentclass[11pt]{article}
\usepackage[margin=1in]{geometry}
\usepackage{amsmath,amssymb,amsthm,hyperref}
\newtheorem{theorem}{Theorem}
\newtheorem{proposition}[theorem]{Proposition}
\newtheorem{lemma}[theorem]{Lemma}
\newtheorem{corollary}[theorem]{Corollary}
\theoremstyle{remark}\newtheorem{remark}[theorem]{Remark}
\newcommand{\E}{\mathbb E}
\newcommand{\Prb}{\mathbb P}
\newcommand{\R}{\mathbb R}
\newcommand{\ind}{\mathbf 1}
\newcommand{\Var}{\operatorname{Var}}
\newcommand{\Cov}{\operatorname{Cov}}
\newcommand{\KL}{\operatorname{KL}}
\newcommand{\TV}{\operatorname{TV}}
\newcommand{\clip}{\operatorname{clip}}
\newcommand{\supp}{\operatorname{supp}}
\DeclareMathOperator*{\argmax}{arg\,max}
\DeclareMathOperator*{\argmin}{arg\,min}
\setlength{\parskip}{5pt}
\setlength{\parindent}{0pt}
\title{Information and statistical power under interference\\[4pt]\large Robust conditional tests and exact designs on paired networks}
\author{GPT 6 Astra Ultra}
\date{9 October 2026}
\begin{document}
\maketitle

\section*{Question and scope}
The target is design-dependent power for spillover effects when exposure is correlated with direct treatment and errors may be dependent. The known Gaussian information formula alone is not a solution. This attempt gives a covariance-robust conditional test, a matching small-information obstruction, and a complete optimal-design calculation on disjoint pairs. It also distinguishes conditional collinearity from unconditional identification under randomized treatment counts.

\section*{Fixed-design robust inference}
There are $N$ vertices with a fixed observed graph. Treatment $w\in\{0,1\}^N$ determines exposure $s=S(w)\in\R^N$. Conditional on this assignment,
\[
 Y=A\theta+\delta s+\varepsilon,\qquad
 A=[X,w],\quad\varepsilon\sim N(0,\Sigma),
\]
where $X$ contains recorded pretreatment regressors, $\theta$ is unrestricted and $\Sigma$ is independent of treatment. Let $\Pi_A$ be the Euclidean orthogonal projection onto the column space of $A$, even if $A$ is rank deficient, and set $r=(I-\Pi_A)s$. The covariance class is
$\mathcal S=\{\Sigma:0<\underline v I\preceq\Sigma\preceq\bar v I\}$ with known $0<\underline v\le\bar v<\infty$. Denote the standard normal distribution function by $\Phi$ and its $p$ quantile by $z_p$.

\begin{theorem}[Uniform conditional size and power]
For $r\ne0$, reject $\delta=0$ when
\[
 |r'Y|>z_{1-\alpha/2}\sqrt{\bar v}\,\|r\|,
 \qquad 0<\alpha<1.
\]
For $r=0$, never reject. The test has size at most $\alpha$ uniformly over $\theta$ and $\Sigma\in\mathcal S$, conditional on every assignment. For $0<\beta<1/2$, its power is at least $1-\beta$ whenever
\[
 |\delta|\|r\|\ge\sqrt{\bar v}
 (z_{1-\alpha/2}+z_{1-\beta}).
\]
For this covariance class the worst-case efficient conditional information is exactly $\|r\|^2/\bar v$.
\end{theorem}
\begin{proof}
Because $r'A=0$ and $r's=\|r\|^2$, the contrast $r'Y$ is normal with mean $\delta\|r\|^2$ and variance $r'\Sigma r\le\bar v\|r\|^2$. A centered normal with smaller variance has smaller two-sided tails at a positive fixed cutoff, proving the size bound. Under a positive mean, the probability of falling below the positive cutoff is at most $\beta$ if the distance from the mean to that cutoff is at least $z_{1-\beta}\sqrt{\bar v}\|r\|$; a negative mean is symmetric. This gives the power bound.

Efficient information is the minimized weighted squared distance
$\min_b(s-Ab)'\Sigma^{-1}(s-Ab)$. Since $\Sigma\preceq\bar v I$, it is at least $\min_b\|s-Ab\|^2/\bar v=\|r\|^2/\bar v$. Equality holds for the admissible covariance $\Sigma=\bar v I$.
\end{proof}

Randomizing $w$ independently of errors preserves size by conditioning and averaging. A uniform conditional power guarantee requires the residual norm bound on every supported assignment; maximizing its expectation alone cannot give such a guarantee.

\begin{theorem}[Exact two-point obstruction]
At a fixed assignment and $\Sigma=\bar v I$, for any $\Delta>0$ there are nuisance values under $\delta=0$ and $\delta=\Delta$ such that the sum of errors of every test is at least
\[
 2\Phi\left(-\frac{\Delta\|r\|}{2\sqrt{\bar v}}\right).
\]
If $r=0$, the two observation laws can be made identical for every $\Delta$, so no test has size $\alpha$ and power greater than $\alpha$ uniformly over nuisance values.
\end{theorem}
\begin{proof}
Choose $b$ with $Ab=\Pi_A s$. Under the null take nuisance $\theta_0=0$; under the alternative take $\theta_1=-\Delta b$. Then the means are zero and $\Delta r$. Their Gaussian likelihood ratio exceeds one exactly when $r'Y>\Delta\|r\|^2/2$. At each observation, assigning it to the larger density minimizes the sum of the two error probabilities, since the error contribution is then the smaller density. For this likelihood-ratio rule both errors equal $\Phi(-\Delta\|r\|/(2\sqrt{\bar v}))$. No other test can improve their sum. If $r=0$, the means and covariances coincide.
\end{proof}

Thus the separation scale $\sqrt{\bar v}/\|r\|$ is unavoidable, up to error-probability constants. It is a feature of the design and exposure relation, not simply graph size or average degree.

\section*{A complete design calculation on disjoint pairs}
Let $N=2n$ and let the graph consist of $n$ disjoint edges. Exposure is the treatment of the unique neighbor. Let $X=\mathbf1$ and impose exactly $n$ treated vertices. Write $m$ for the number of mixed edges (one treated endpoint). The other edges are equally divided between types $00$ and $11$, so $0\le m\le n$ and $n-m$ is even.

\begin{theorem}[Paired-network optimality]
For every such assignment,
\[
 \|(I-\Pi_{[\mathbf1,w]})S(w)\|^2
 =2m\left(1-\frac mn\right).
\]
Consequently the assignments maximizing worst-case conditional information are exactly those whose feasible integer $m$ is closest to $n/2$. If $N$ is a multiple of eight, the optimum is $m=n/2$ and the residual norm squared is $N/4$. Both pure cluster assignment ($m=0$) and exactly one treated vertex per edge ($m=n$) have zero information. A randomized design supported entirely on the optimal assignments achieves the same information guarantee while retaining random treatment assignment.
\end{theorem}
\begin{proof}
Both $w$ and $s=S(w)$ have exactly $n$ ones, so their centered versions have squared norm $N/4=n/2$. Each mixed edge contributes $-1/2$ to their inner product and each unmixed edge contributes $1/2$. Thus their inner product is $(n-2m)/2$, and their correlation is $1-2m/n$. Regressing centered $s$ on centered $w$ leaves squared residual
\[
 \frac n2\left[1-\left(1-\frac{2m}n\right)^2\right]
 =2m(1-m/n).
\]
This concave quadratic is maximized by feasible integers nearest its vertex $n/2$. When $N=8k$, $n=4k$ and $m=2k$ is feasible because $n-m=2k$ is even. Substitution gives $n/2=N/4$. The zero-information endpoints follow directly. Since all assignments in the proposed randomization have the same residual norm, the conditional information and test guarantees hold for every draw.
\end{proof}

A concrete optimal design selects $n/4$ edges to be $00$, $n/4$ to be $11$ and $n/2$ to be mixed, randomizing these labels uniformly and the treated endpoint of every mixed edge fairly. Every vertex then has marginal treatment probability $1/2$. It satisfies both the exact budget and the information optimum; correlation across assignments is part of the design, not a violation of randomization.

\section*{Complete graphs and the random-count qualification}
\begin{proposition}[A global nonidentification example]
On the complete graph with $N\ge2$, define neighbor-average exposure
$S_i(w)=(\sum_jw_j-w_i)/(N-1)$. With an intercept and exactly $k$ treated vertices in every supported assignment, the joint law of $(W,Y)$ depends on $(\gamma,\beta,\delta)$ only through
\[
 \gamma+\delta k/(N-1),\qquad \beta-\delta/(N-1),
\]
where $\gamma$ is the intercept and $\beta$ the direct effect. Thus $\delta$ is not identified even unconditionally. If the design instead supports two different treatment counts and contains nonconstant assignments at those counts, conditional collinearity by itself no longer proves this nonidentification.
\end{proposition}
\begin{proof}
For fixed count $k$, $S(w)=k\mathbf1/(N-1)-w/(N-1)$ with the same coefficients for every assignment. Substitution gives the stated two combinations, so changing $\delta$ and offsetting $\gamma,\beta$ produces identical joint laws.

For varying counts, suppose two parameter triples have identical conditional means for every supported assignment. On a nonconstant assignment, the coefficients on $\mathbf1$ and $w$ must separately agree, giving $\Delta\beta=\Delta\delta/(N-1)$ and $\Delta\gamma=-\Delta\delta k/(N-1)$. Two different counts force $\Delta\delta=0$. Thus in this support configuration varying treatment counts removes this mean-level aliasing, even though each individual assignment has zero residual exposure after fitting its own intercept and direct effect.
\end{proof}

\section*{Independent network replicates and the remaining problem}
If the same realized assignment (and therefore the same contrast $r$) is implemented on $M$ independent networks with the same $(\theta,\delta,\Sigma)$, the scalar contrasts $r'Y_j/\|r\|^2$ are independent normal observations with mean $\delta$ and variance $r'\Sigma r/\|r\|^4$. Their average has variance divided by $M$, giving the preceding known-bound test with separation divided by $\sqrt M$. This conclusion follows by multiplying the normal moment-generating functions. Dependence between those networks would invalidate that replicate count; independence must concern the errors and the intervention system, not merely labels assigned to graph pieces.

For general graphs, finding the best residual geometry subject to treatment constraints is a nontrivial combinatorial problem. The spectral covariance bound used here is deliberately conservative and does not solve power optimization under rich estimated covariance classes. Optimal unconditional tests with fixed nuisance across assignments can exploit information discarded by conditional projection, as the complete-graph example shows. These general questions remain unresolved.

The primary survey text's Section 3.2.3 was checked for its statistical-power and replicates discussion. The paired-graph optimum and covariance-class specialization above are deductions for this explicitly stated model, not claimed solutions attributed to the survey.
\begin{thebibliography}{9}
\bibitem{c} C. Cinelli, A. Feller, G. Imbens, E. Kennedy, S. Magliacane and J. Zubizarreta (2025), Challenges in Statistics: A Dozen Challenges in Causality and Causal Inference, arXiv:2508.17099v1, Section 3.2.3. \url{https://arxiv.org/html/2508.17099v1}.
\end{thebibliography}

\end{document}
